\section{Estado del arte}

El panorama actual del descubrimiento de estructuras en redes bayesianas y modelos gráficos causales ha experimentado una notable evolución, transitando desde la formulación clásica de algoritmos basados en restricciones hacia métodos altamente escalables, flexibles e híbridos. Para comprender de manera sistemática este estado del arte, la presente sección se organiza en tres ejes analíticos interconectados. En primer lugar, se examina la evolución del algoritmo PC y sus variantes, abordando desde sus fundamentos teóricos y modificaciones para garantizar la independencia del orden de entrada y la paralelización, hasta sus extensiones modernas para datos temporales, valores extremos y optimización diferenciable. En segundo lugar, se analizan las pruebas de hipótesis de independencia condicional, comprendidas como el motor estadístico crítico que sustenta las decisiones de poda de aristas y orientación en presencia de variados tipos de datos. Por último, se exploran los paradigmas híbridos, los cuales integran el filtrado local por restricciones con la optimización por puntuación o la resolución lógica de conflictos para mitigar la acumulación de errores en muestras finitas y mejorar la precisión estructural.

\subsection{Evolución del Algoritmo PC: Fundamentos, Inestabilidad y Variantes Algorítmicas}

El \textbf{algoritmo PC} (denominado así por sus creadores Peter Spirtes y Clark Glymour) constituye el método clásico y fundacional dentro del paradigma de aprendizaje de estructuras causales basado en restricciones. Desarrollado como una alternativa computacionalmente eficiente al algoritmo SGS (\textit{Spirtes, Glymour y Scheines}), el algoritmo PC parte de un grafo completo no dirigido y elimina aristas de forma recursiva basándose en decisiones de independencia condicional. Esto produce un grafo no dirigido que luego puede ser parcialmente orientado y extendido para representar el DAG subyacente. El algoritmo PC se ejecuta en el peor de los casos en tiempo exponencial (en función del número de nodos), pero si el DAG subyacente real es disperso, esta complejidad se reduce a un tiempo de ejecución polinomial. \citep{spirtes2000causation, scanagatta2019survey}

En el algoritmo PC, asumimos un conjunto de variables $\mathbf{X} = (X_1, \ldots, X_n)$ con una distribución de probabilidad global $P$ sobre ellas. Los subconjuntos de variables de $\mathbf{X}$ se representan mediante letras mayúsculas en negrita (por ejemplo, $\mathbf{A}, \mathbf{B}, \mathbf{C}$). La notación $I(\mathbf{A}, \mathbf{B} \mid \mathbf{C})$ denota que los conjuntos de variables $\mathbf{A}$ y $\mathbf{B}$ son condicionalmente independientes dado el conjunto $\mathbf{C}$. El algoritmo PC asume la hipótesis de fidelidad (\textit{faithfulness}). Esto significa que existe un grafo acíclico dirigido (DAG), $G$, tal que las relaciones de independencia entre las variables de $\mathbf{X}$ son exactamente aquellas representadas por $G$ mediante el criterio de $d$-separación. En consecuencia, el algoritmo PC se apoya en la posibilidad de comprobar cuándo se verifica $I(\mathbf{A}, \mathbf{B} \mid \mathbf{C})$ en el grafo $G$ y, a partir de esa información, recuperar la estructura subyacente. \citep{spirtes2000causation,abellan2006variations}

La idea central del método consiste en recuperar primero el esqueleto del grafo y, posteriormente, orientar las aristas que permanecen. La etapa de poda se realiza eliminando enlaces entre parejas de nodos cuando existe un subconjunto de sus vecinos que hace que ambas variables sean condicionalmente independientes dados esos vecinos. Este criterio permite reducir progresivamente el grafo completo inicial hasta un grafo no dirigido compatible con las independencias observadas. \citep{spirtes2000causation, abellan2006variations}


\begin{algorithm}[H]
\caption{Algoritmo PC: fase de búsqueda del esqueleto}
\begin{algorithmic}[1]
\State Inicializar $G'$ como un grafo completo no dirigido sobre $\mathbf{X}$
\State $i \gets 0$
\Repeat
    \For{cada par adyacente $(X,Y)$ en $G'$}
        \If{existe $S \subseteq \mathrm{Adj}_{G'}(X) \setminus \{Y\}$ con $|S| = i$ y $I(X,Y \mid S)$}
            \State Guardar $S_{XY} \gets S$
            \State Eliminar la arista $X-Y$ de $G'$
        \EndIf
    \EndFor
    \State $i \gets i + 1$
\Until{$\forall X \in \mathbf{X},\; |\mathrm{Adj}_{G'}(X)| \le i$}
\end{algorithmic}
\end{algorithm}

En que el conjunto completo de variables del problema se denota por $\mathbf{X}$, mientras que $\mathrm{Adj}_{G'}(X)$ representa el conjunto de vecinos de $X$ en el grafo actual $G'$. La expresión $I(X,Y \mid S)$ indica que $X$ e $Y$ son condicionalmente independientes dado el conjunto $S$. Así, si las independencias observadas son fieles a un DAG, no existe un enlace entre $X$ e $Y$ si y solo si existe un subconjunto $S$ de los vecinos de $X$ tal que se cumple $I(X,Y \mid S)$. Para cada par de variables, $S_{XY}$ contendrá dicho conjunto, si se encuentra. Este conjunto se utilizará en la etapa posterior de orientación.\citep{spirtes2000causation, abellan2006variations}

La etapa de orientación procederá buscando conjuntos de tres variables $\{X, Y, Z\}$ tales que los enlaces $X - Z$ y $Y - Z$ estén en el grafo pero no el enlace $X - Y$. Entonces, si $Z \notin S_{XY}$, orienta los enlaces de $X$ a $Z$ y de $Y$ a $Z$ creando una estructura en v: $X \rightarrow Z \leftarrow Y$. Una vez realizadas estas orientaciones, se intenta orientar el resto de las aristas siguiendo dos principios básicos: no crear ciclos y no crear nuevas estructuras en v. Es posible que la orientación de algunas aristas deba elegirse arbitrariamente. Si el conjunto de independencias es fiel a un grafo y disponemos de un procedimiento perfecto para determinar si $I(X, Y \mid S)$, entonces el algoritmo garantiza producir un grafo equivalente (que representa el mismo conjunto de independencias) al original. \citep{abellan2006variations}

A pesar de su solidez teórica y su completitud asintótica bajo supuestos de fidelidad y suficiencia causal, el algoritmo PC clásico presenta varias limitaciones críticas cuando se aplica a datos reales de tamaño muestral finito \citep{colombo2014order}. En primer lugar, la versión muestral del algoritmo sufre una fuerte dependencia del orden, donde la secuencia en la que se disponen las variables en el conjunto de datos determina qué pruebas de independencia condicional se ejecutan primero; como resultado, los errores estadísticos tempranos alteran dinámicamente los conjuntos de condicionamiento disponibles para pares de variables posteriores, lo que conduce a esqueletos estimados, estructuras en $v$ y orientaciones de aristas altamente variables. En segundo lugar, las pruebas estadísticas frecuentistas estándar (como $\chi^2$ o las pruebas $z$ de correlación parcial) toman decisiones categóricas e irrevocables desde el principio, y se vuelven cada vez más propensas a errores de tipo I y tipo II a medida que disminuyen los tamaños de muestra o crece la cardinalidad del conjunto de condicionamiento \citep{abellan2006variations}. En tercer lugar, aunque el algoritmo es computacionalmente polinómico para grafos dispersos, el cálculo de numerosas pruebas de independencia condicional en conjuntos de datos de alta dimensión o con muestras masivas crea un cuello de botella computacional dominante. Estos desafíos combinados —inestabilidad del orden, vulnerabilidad estadística en muestras finitas y limitaciones de escalabilidad computacional— motivaron el desarrollo de las variantes algorítmicas especializadas que se discuten a continuación.

Abellán, Gómez-Olmedo y Moral (2006) propusieron cuatro modificaciones clave al algoritmo PC estándar para abordar las limitaciones causadas por errores estadísticos en contextos de muestras finitas. Para minimizar los errores de decisión asociados a conjuntos de condicionamiento grandes, introdujeron la evaluación de la independencia condicional mediante conjuntos de corte de tamaño mínimo ($CUT_{X,Y}$), calculados con un algoritmo de Ford-Fulkerson modificado. Además, sustituyeron las pruebas clásicas de $\chi^2$ por una métrica de puntuación bayesiana de Dirichlet (BD), que ofrece un equilibrio más razonable entre los errores de tipo I y tipo II en distintos tamaños de muestra. Asimismo, incorporaron un paso de refinamiento greedy posterior para optimizar el grafo aprendido según una puntuación BD con el fin de aproximar mejor la distribución conjunta de probabilidad, y un método de resolución de triángulos que evalúa la fuerza de los enlaces para podar el arco más débil en subgrafos totalmente conectados de tres variables.

Kalisch y Bühlmann (2007) establecieron la validez teórica de aplicar el algoritmo PC en escenarios de alta dimensión, donde el número de variables $p$ crece rápidamente en función del tamaño muestral $n$ ($p = O(n^a)$ con $0 < a < \infty$). Demostraron que la versión muestral del algoritmo PC, utilizando la transformación $z$ de Fisher para pruebas de correlación parcial, es asintóticamente consistente para estimar el esqueleto de DAG dispersos bajo distribuciones gaussianas multivariantes. Bajo la suposición de que los vecindarios del grafo son dispersos ($q_n = O(n^{1-b})$ con $0 < b \le 1$), mostraron que la complejidad computacional en el peor caso queda acotada por $O(p^{q_n})$, lo que hace que el aprendizaje de estructuras sea computacionalmente tratable en redes altas y dispersas. Por último, sus simulaciones empíricas demostraron que el algoritmo PC alcanza tasas de descubrimiento verdaderas (TDR) significativamente mayores que las alternativas basadas en puntuación y mantiene una gran robustez en un amplio rango de niveles de significación $\alpha$.

Kalisch et al. (2012) presentaron el paquete de código abierto para R \texttt{pcalg}, diseñado para facilitar el aprendizaje de estructuras causales y la estimación cuantitativa de efectos de intervenciones causales a partir de datos observacionales. El paquete implementa varios algoritmos fundamentales de descubrimiento de estructuras basados en restricciones: el algoritmo PC para escenarios con suficiencia causal, y los algoritmos FCI y RFCI, que contemplan la presencia de factores de confusión latentes y variables de selección. Además de estimar estructuras, \texttt{pcalg} incorpora el método IDA y su variante optimizada desde el punto de vista computacional, \texttt{idaFast()}, que aplican el cálculo de intervenciones de Pearl a clases de equivalencia de Markov aprendidas para estimar cotas de los efectos causales totales cuando el grafo subyacente no puede identificarse de forma única. Para adaptarse a diversos contextos empíricos, el paquete ofrece funciones modulares para realizar pruebas de independencia condicional con datos gaussianos, discretos y binarios, y permite integrar procedimientos de prueba personalizados y conocimiento previo, como aristas obligatorias o prohibidas. También incluye opciones como reglas conservadoras de orientación para mitigar los conflictos entre aristas inducidos por errores de muestreo. Su utilidad práctica se ha demostrado en aplicaciones de biología y epidemiología de alta dimensionalidad, con miles de variables \citep{kalisch2012causal}.

Colombo y Maathuis (2014) abordaron la cuestión crítica de la dependencia del orden en los algoritmos de aprendizaje de estructuras causales basados en restricciones (incluidos PC, FCI, RFCI y CCD), demostrando que, en escenarios de alta dimensión y muestras finitas, la variación en la secuencia de variables de entrada produce una inestabilidad sustancial en los esqueletos estimados, las estructuras en v y las orientaciones de aristas. Para eliminar la dependencia del orden en la búsqueda del esqueleto, introdujeron el algoritmo \textbf{PC-stable}, que registra y congela los conjuntos de adyacencia $a(X_i)$ al inicio de cada nivel de condicionamiento $\ell$, impidiendo que las eliminaciones de aristas dentro del nivel $\ell$ alteren prematuramente los conjuntos de condicionamiento para otros pares de variables. Para la determinación de estructuras en v independientes del orden, desarrollaron \textbf{CPC-stable} (Conservative PC) y \textbf{MPC-stable} (Majority Rule PC), que evalúan todos los conjuntos separadores para triples no protegidos y clasifican un triple como no ambiguo solo si el nodo intermedio es consistente en su presencia o ausencia en todos los conjuntos separadores (CPC) o en un umbral mayoritario de ellos (MPC). Además, propusieron reglas de orientación basadas en listas (\textbf{LCPC-stable} y \textbf{LMPC-stable}) que utilizan aristas bidireccionales para resolver conflictos de orientación, logrando salidas gráficas completamente independientes del orden. De manera crucial, demostraron que todas las modificaciones propuestas preservan la corrección, la completitud y la consistencia en alta dimensión del algoritmo PC original bajo condiciones teóricas idénticas, mientras extienden con éxito estos principios de independencia del orden a algoritmos que manejan variables latentes como FCI y RFCI. A través de simulaciones de alta dimensión y análisis de expresión genética, demostraron que las variantes PC-stable producen una varianza significativamente menor de la Distancia de Hamming Estructural (SHD), tasas de descubrimiento verdaderas (TDR) superiores y predicciones de efectos causales mejores cuando se integran con IDA.

Le et al. (2019) abordaron el importante cuello de botella computacional que supone el tiempo de ejecución de los métodos de descubrimiento causal basados en restricciones e independientes del orden, especialmente en aplicaciones biológicas de alta dimensionalidad. Aunque variantes como PC-stable eliminan la dependencia del orden al fijar las consultas de vecindad en cada nivel de condicionamiento, requieren evaluar un número considerablemente mayor de pruebas de independencia condicional, lo que puede producir tiempos de ejecución prohibitivos —superiores a 24 horas— en grandes conjuntos de datos de expresión génica. Para reducir este coste sin sacrificar la independencia del orden ni la precisión estructural, desarrollaron \textbf{Parallel-PC}, un marco de paralelización eficiente en memoria para equipos de escritorio multinúcleo y servidores de laboratorio. Como las pruebas de independencia condicional de un mismo nivel son mutuamente independientes mientras las adyacencias del grafo permanecen fijas, Parallel-PC distribuye esas pruebas entre los núcleos disponibles y sincroniza las actualizaciones al pasar de un nivel al siguiente. Además, asigna a un mismo núcleo las pruebas asociadas a una arista candidata, evitando evaluaciones redundantes entre conjuntos de nodos adyacentes. En pruebas de alta dimensionalidad, incluidos los conjuntos de datos del desafío DREAM5, redujo el tiempo de ejecución de más de 24 horas a menos de seis horas en una máquina de ocho núcleos. Su integración en el flujo de trabajo \texttt{efficient-IDA} para identificar dianas reguladoras de miARN y ARNm también mejoró la escalabilidad computacional y la precisión de la predicción de efectos causales en redes biológicas complejas \citep{le2016fast}.

Kocaoglu (2023) abordó la limitación del algoritmo PC y sus derivados en escenarios con un volumen de datos restringido, en los que las pruebas de independencia condicional pierden potencia estadística a medida que aumenta el tamaño del conjunto condicionante. Para evitar la propagación en cadena de errores estadísticos generados por pruebas de alto orden, el autor introdujo la noción de \textbf{$k$-equivalencia de Markov}, según la cual dos grafos acíclicos dirigidos (DAGs) se consideran equivalentes si comparten las mismas restricciones de independencia condicional para conjuntos condicionantes de tamaño limitado a un entero $k$ ($|C| \leq k$). Para representar esta nueva clase de equivalencia, el estudio definió los \textbf{grafos de $k$-clausura} (\textit{$k$-closure graphs}) y demostró teóricamente que la $k$-equivalencia entre dos DAGs se reduce a la equivalencia de Markov de sus grafos de $k$-clausura. Sobre esta caracterización, el autor desarrolló el algoritmo \textbf{$k$-PC}, que adapta las reglas de orientación del algoritmo FCI para aprender la estructura del grafo $k$-esencial restringiendo las evaluaciones a pruebas de orden como máximo $k$. Experimentos con datos sintéticos y con la red \textit{Asia} demostraron que el algoritmo $k$-PC presenta una mayor robustez en regímenes de muestra reducida, alcanzando puntuaciones F1 significativamente superiores en la identificación del esqueleto y la orientación de las aristas en comparación con el PC tradicional y métodos de optimización continua como \textbf{NOTEARS}.

Biswas y Mukherjee (2024) ampliaron los fundamentos matemáticos del algoritmo PC a escenarios con datos no independientes e idénticamente distribuidos, y establecieron garantías teóricas de consistencia para el descubrimiento causal en series temporales multivariantes. Mientras que las pruebas clásicas de consistencia uniforme del algoritmo PC suponen muestras independientes e idénticamente distribuidas, las series temporales presentan dependencia temporal de forma natural. Los autores demostraron que el algoritmo PC muestral conserva la consistencia puntual y asintótica al estimar DAG contemporáneos bajo procesos estacionarios con dependencia de mezcla $\rho$ y de mezcla fuerte, siempre que las correlaciones decaigan con suficiente rapidez. Establecieron estas garantías para estructuras estacionarias habituales, incluidos los modelos de medias móviles autorregresivas vectoriales (VARMA) y los procesos lineales, utilizando tanto pruebas paramétricas $z$ de correlación parcial de Pearson como pruebas no paramétricas basadas en el criterio de independencia de Hilbert-Schmidt (HSIC). Asimismo, demostraron la consistencia asintótica en alta dimensionalidad del algoritmo \textbf{Time-Aware PC (TPC)}, que aprende grafos causales entre instantes y rezagos temporales y luego los transforma en grafos causales resumidos y desplegados. En experimentos con modelos VAR gaussianos lineales y no lineales, redes neuronales de tiempo continuo y datos hidrológicos reales del conjunto CauseMe River Runoff, TPC superó al algoritmo PC clásico y a la causalidad de Granger en la recuperación de conexiones causales entre instantes \citep{biswas2024consistent}.

Las principales variantes y extensiones del algoritmo PC analizadas en esta subsección se resumen en la Tabla~\ref{tab:variantes_pc}, incluida en el Apéndice~\ref{app:variantes_pc}.

\subsection{Pruebas de Independencia Condicional y Extensiones para Causalidad}

La base del aprendizaje estructural de tipo \emph{constraint-based} reside en las pruebas de independencia condicional, que permiten decidir si dos variables son dependientes o independientes dadas otras variables del sistema...

\subsection{Enfoques Híbridos (Constraint-Based y Score-and-Search) y Líneas Futuras}

La literatura reciente ha mostrado un creciente interés en enfoques híbridos que combinan la fuerza de los métodos\emph{constraint-based} con las estrategias de\emph{score-and-search}. ...
